% (c) Nikita Lisitsa, lisyarus@gmail.com, 2026

\documentclass[10pt]{beamer}

\usepackage[T2A]{fontenc}
\usepackage[russian]{babel}
\usepackage{minted}

\usepackage[most]{tcolorbox}
\tcbuselibrary{minted}

\usepackage{graphicx}
\graphicspath{ {./images/} }

\usepackage{adjustbox}

\usepackage{color}
\usepackage{soul}
\usepackage{multicol}
\usepackage{hyperref}
\usepackage{tikz}
\usetikzlibrary{arrows.meta,positioning,calc}

\usetheme{metropolis}

\definecolor{red}{rgb}{1,0,0}
\definecolor{codebg}{RGB}{29,35,49}
\setminted{fontsize=\footnotesize}

\makeatletter
\newcommand{\slideimage}[1]{
  \begin{figure}
    \begin{adjustbox}{width=\textwidth, totalheight=\textheight-2\baselineskip-2\baselineskip,keepaspectratio}
      \includegraphics{#1}
    \end{adjustbox}
  \end{figure}
}
\makeatother

\title{Язык программирования C++}
\subtitle{Лекция 16: STL, стандартные контейнеры, итераторы, инвалидация ссылок, инвалидация итераторов}
\date{2026}
\author{Никита Лисица}

\setbeamertemplate{footline}[frame number]

\begin{document}

\frame{\titlepage}

\begin{frame}[fragile]
\frametitle{Напоминание: стандартная библиотека}
\begin{itemize}
\usemintedstyle{tango}
\item Поставляется вместе с реализацией (компилятором) языка
\pause
\item Заголовочные файлы не имеют расширения:
\usemintedstyle{lightbulb}
\begin{tcblisting}{listing engine=minted, minted language=cpp, minted options={fontsize=\scriptsize}, listing only, colback=codebg, colframe=codebg, rounded corners}
#include <vector>
#include <string>
\end{tcblisting}
\pause
\usemintedstyle{tango}
\item Большинство функций и классов лежат в \mintinline{cpp}|namespace std|:
\usemintedstyle{lightbulb}
\begin{tcblisting}{listing engine=minted, minted language=cpp, minted options={fontsize=\scriptsize}, listing only, colback=codebg, colframe=codebg, rounded corners}
std::vector<int> v = {1, 2, 3, 4};
std::string str = "I love Rust";
\end{tcblisting}
\pause
\item Широко использует всё, о чём мы говорили ранее в курсе: классы, шаблоны, исключения
\end{itemize}
\end{frame}

\begin{frame}[fragile]
\frametitle{STL}
\begin{itemize}
\usemintedstyle{tango}
\item Standard Template Library (STL) -- часть стандартной библиотеки, предоставляющая
\pause
\begin{itemize}
\item Контейнеры (структуры данных для хранения и поиска)
\pause
\item Алгоритмы (для обработки данных)
\end{itemize}
\pause
\item \textbf{NB}: часто термином STL называют всю стандартную библиотеку, что неверно
\end{itemize}
\end{frame}

\begin{frame}[fragile]
\frametitle{Контейнеры}
\begin{itemize}
\usemintedstyle{tango}
\item Шаблоны классов, реализующие различные структуры данных
\pause
\item Параметризуются типом хранимых объектов, от которых требуется
\begin{itemize}
\item Иметь конструктор копирования (или move-конструктор)
\item Иметь оператор присваивания (или move-присваивания)
\item Иметь деструктор
\end{itemize}
\end{itemize}
\end{frame}

\begin{frame}[fragile]
\frametitle{Контейнеры}
\begin{itemize}
\usemintedstyle{tango}
\item Все контейнеры имеют
\pause
\begin{itemize}
\item Конструктор по умолчанию
\pause
\item Конструктор копирования и оператор присваивания
\pause
\item Деструктор
\pause
\item \mintinline{cpp}|size()| -- размер контейнера
\pause
\item \mintinline{cpp}|empty()| -- проверить на пустоту
\pause
\item \mintinline{cpp}|clear()| -- очистить (удалить все элементы)
\pause
\item \mintinline{cpp}|swap(other)|
\pause
\item \mintinline{cpp}|insert(...)|
\pause
\item \mintinline{cpp}|erase(...)|
\end{itemize}
\end{itemize}
\end{frame}

\begin{frame}[fragile]
\frametitle{Последовательые контейнеры}
\begin{itemize}
\usemintedstyle{tango}
\item Сохраняют порядок, в котором были добавлены элементы:
\pause
\begin{itemize}
\item \mintinline{cpp}|std::vector| -- саморасширяющийся массив
\pause
\item \mintinline{cpp}|std::list| -- двусвязный список
\pause
\item \mintinline{cpp}|std::deque| -- дек (double-ended queue)
\pause
\item \mintinline{cpp}|std::string| -- строка (массив символов)
\pause
\item \mintinline{cpp}|std::array| -- массив фиксированного размера (аналог \mintinline{cpp}|T[N]|)
\pause
\item \mintinline{cpp}|std::forward_list| -- односвязный список
\end{itemize}
\end{itemize}
\end{frame}

\begin{frame}[fragile]
\frametitle{std::vector}
\begin{itemize}
\usemintedstyle{tango}
\item Саморасширяющийся массив: хранит элементы непрерывно в памяти, выделенной с запасом под новые элементы
\pause
\item При исчерпании запаса, увеличивает размер (часто -- в 2 раза) реаллоцируя память и копируя элементы
\pause
\item \mintinline{cpp}|size()| -- количество элементов, \mintinline{cpp}|capacity()| -- общий размер памяти (в элементах, не в байтах)
\end{itemize}
\end{frame}

\begin{frame}[fragile]
\frametitle{std::vector}
\begin{itemize}
\usemintedstyle{tango}
\item Добавление/удаление в конец -- амортизированное \(O(1)\)
\pause
\item Добавление/удаление в середине -- \(O(N)\), так как нужно сдвинуть остальные элементы
\pause
\item Доступ к элементу по индексу -- \(O(1)\)
\pause
\item Из-за кешей CPU, prefetching'а и т.п., \mintinline{cpp}|vector| -- часто самый эффективный из контейнеров, и почти всегда это хороший дефолт
\end{itemize}
\end{frame}

\begin{frame}[fragile]
\frametitle{std::vector}
\begin{itemize}
\usemintedstyle{tango}
\item \mintinline{cpp}|resize(n)| -- изменить размер (количество реально хранимых элементов)
\pause
\item \mintinline{cpp}|reserve(n)| -- изменить размер памяти (запаса)
\pause
\item \mintinline{cpp}|push_back(x)/pop_back()| -- добавить/удалить последний элемент
\pause
\item \mintinline{cpp}|front()/back()| -- получить ссылку на первый/последний элемент
\pause
\item \mintinline{cpp}|operator[]| -- доступ по индексу (выход за границу -- UB)
\pause
\item \mintinline{cpp}|at(i)| -- доступ по индексу (выход за границу -- исключение)
\pause
\item \mintinline{cpp}|data()| -- получить указатель на хранимый массив (напр. для взаимодействия с библиотеками на C)
\end{itemize}
\end{frame}

\begin{frame}[fragile]
\frametitle{std::vector}
\begin{itemize}
\usemintedstyle{lightbulb}
\begin{tcblisting}{listing engine=minted, minted language=cpp, minted options={fontsize=\scriptsize}, listing only, colback=codebg, colframe=codebg, rounded corners}
#include <vector>

std::vector<int> v(10);
v.push_back(42); // 11ый элемент
v[0] = 13;
v.at(5) = 14;
// v.at(11) -- бросит исключение
v.size(); // вернёт 11
v.pop_back();
v.empty(); // вернёт false
v.clear(); // после этого v.empty() == true

// Передаём данные в сишный API
glBufferData(v.data(), ...);
\end{tcblisting}
\end{itemize}
\end{frame}

\begin{frame}[fragile]
\frametitle{std::vector}
\begin{itemize}
\usemintedstyle{lightbulb}
\begin{tcblisting}{listing engine=minted, minted language=cpp, minted options={fontsize=\scriptsize}, listing only, colback=codebg, colframe=codebg, rounded corners}
std::vector<int> v(10);

// Удалит все элементы кроме первых 5
v.resize(5);
// Добавит недостающие эелементы конструктором по умолчанию
v.resize(20);
// Можно сказать, чем нужно заполнить новые элементы
v.resize(30, 5);

// Выделим память под 100 элементов (размер не изменится)
v.reserve(100);
// Удалим все элементы, но запас памяти останется под новые
v.clear();

// Загадка: что это и зачем?
std::vector<int>().swap(v);
// То же самое:
v.clear(); v.shrink_to_fit();
\end{tcblisting}
\end{itemize}
\end{frame}

\begin{frame}[fragile]
\frametitle{std::deque}
\begin{itemize}
\usemintedstyle{tango}
\item Double-ended queue: добавление/удаление с обоих концов за \(O(1)\)
\pause
\item Обычно реализуется как массив указателей на массивы фиксированного размера (или даже вектор векторов)
\pause
\item Гарантирует постоянство ссылок на элементы (иначе его можно было бы реализовать как циклический массив)
\pause
\item Обращение по индексу за \(O(1)\)
\end{itemize}
\end{frame}

\begin{frame}[fragile]
\frametitle{std::deque}
\begin{itemize}
\usemintedstyle{tango}
\item \mintinline{cpp}|size()/resize(n)|
\pause
\item \mintinline{cpp}|push_back(x)/pop_back()|
\pause
\item \mintinline{cpp}|push_front(x)/pop_front()|
\pause
\item \mintinline{cpp}|front()/back()|
\pause
\item \mintinline{cpp}|operator[]/at(i)|
\end{itemize}
\end{frame}

\begin{frame}[fragile]
\frametitle{std::list}
\begin{itemize}
\usemintedstyle{tango}
\item Обычный двусвязный список
\pause
\item Для удобства реализации на концах всегда есть вспомогательные вершины, не хранящие данные
\pause
\item В качестве оптимизации эти вершины склеены в одну (т.е. список циклический) и хранятся полем в самом списке
\pause
\item Вставка/удаление в любом месте за \(O(1)\)
\pause
\item Нет обращения по индексу
\end{itemize}
\end{frame}

\begin{frame}[fragile]
\frametitle{std::list}
\begin{itemize}
\usemintedstyle{tango}
\item \mintinline{cpp}|size()/resize(n)|
\pause
\item \mintinline{cpp}|push_back(x)/pop_back()|
\pause
\item \mintinline{cpp}|push_front(x)/pop_front()|
\pause
\item \mintinline{cpp}|merge(list)| -- объединить два сортированных списка в один
\pause
\item \mintinline{cpp}|splice(...)| -- переместить часть одного списка в другой список
\end{itemize}
\end{frame}

\begin{frame}[fragile]
\frametitle{Драма о std::list::splice}
\begin{itemize}
\usemintedstyle{tango}
\item По сути, \mintinline{cpp}|splice()| сводится к переподвешиванию указателей в вершинах списка, и должен выполняться за \(O(1)\)
\pause
\item Проблема в том, как вычислять размер списка
\pause
\item За \(O(1)\) -- хранить поле \mintinline{cpp}|size| и обновлять во всех операциях, тогда \mintinline{cpp}|splice()| придётся за \(O(N)\) вычислять, сколько элементов было перенесено из одного списка в другой
\pause
\item За \(O(N)\) -- вычислять проходом по списку, тогда \mintinline{cpp}|splice()| может работать за \(O(1)\)
\pause
\item До C++11 не было специфицировано, какой именно вариант предоставляет стандартная библиотека
\pause
\item После C++11: \mintinline{cpp}|size()| за \(O(1)\), \mintinline{cpp}|splice()| за \(O(N)\)
\end{itemize}
\end{frame}

\begin{frame}[fragile]
\frametitle{std::string}
\begin{itemize}
\usemintedstyle{tango}
\item Как \mintinline{cpp}|std::vector<char>|, но с особенностями специфичными для строк
\pause
\item Для совместимости с C в конце всегда хранит нулевой символ (не учитывается в \mintinline{cpp}|size()|)
\pause
\item Есть неявный конструктор от строковых литералов \mintinline{cpp}|string(const char *)|
\pause
\item Есть конкатенация строк через \mintinline{cpp}|operator+|
\pause
\item Множество специфичных для строк алгоритмов: \mintinline{cpp}|find()|, \mintinline{cpp}|substr()|, и т.п.
\end{itemize}
\end{frame}

\begin{frame}[fragile]
\frametitle{std::string}
\begin{itemize}
\usemintedstyle{lightbulb}
\begin{tcblisting}{listing engine=minted, minted language=cpp, minted options={fontsize=\scriptsize}, listing only, colback=codebg, colframe=codebg, rounded corners}
std::string str = "Hello";
str += ", ";
std::string str2 = "world!";
std::string str3 = str + str2;

if (str3.find(',') != std::string::npos) {
  ...
}
\end{tcblisting}
\end{itemize}
\end{frame}

\begin{frame}[fragile]
\frametitle{std::string}
\begin{itemize}
\usemintedstyle{tango}
\item Сам \mintinline{cpp}|std::string| -- не шаблон, но специализация шаблона \mintinline{cpp}|std::basic_string<char>|
\pause
\item \mintinline{cpp}|std::wstring = std::basic_string<wchar_t>| -- строка \textit{широких символов} (UTF-32 под Linux, UTF-16 под Windows, часто используются в WinAPI)
\pause
\item \mintinline{cpp}|std::u8string = std::basic_string<char8_t>| -- всегда UTF-8
\pause
\item \mintinline{cpp}|std::u16string = std::basic_string<char16_t>| -- всегда UTF-16
\pause
\item \mintinline{cpp}|std::u32string = std::basic_string<char32_t>| -- всегда UTF-32
\end{itemize}
\end{frame}

\begin{frame}[fragile]
\frametitle{Адаптеры контейнеров}
\begin{itemize}
\usemintedstyle{tango}
\item Реализуют интерфейс некоторого контейнера
\pause
\item Параметризуются другим контейнером, используемым для хранения элементов
\pause
\item \mintinline{cpp}|std::stack<int, std::vector<int>>| -- стек (есть \mintinline{cpp}|push()/pop()|), использующий вектор для хранения
\pause
\item \mintinline{cpp}|std::queue<int, std::list<int>>| -- очередь (есть \mintinline{cpp}|push()/pop()|), использующая список для хранения
\pause
\item \mintinline{cpp}|std::priority_queue<int, std::deque<int>>| -- приоритетная очередь (куча), использующая деку для хранения
\pause
\item \textbf{NB}: для работы с приоритетной очередью часто удобнее использовать свободные функции \mintinline{cpp}|std::make_heap/push_heap/pop_heap|
\end{itemize}
\end{frame}

\begin{frame}[fragile]
\frametitle{Итераторы}
\begin{itemize}
\usemintedstyle{tango}
\item В общем случае -- объект, позволяющий итерироваться по последовательности элементов некоторой коллекции/контейнера
\pause
\item Есть два стиля реализации итераторов -- условно, Java и C++
\end{itemize}
\end{frame}

\begin{frame}[fragile]
\frametitle{Итераторы: Java style}
\begin{itemize}
\usemintedstyle{tango}
\item Java-style: итератор позволяет один раз пройтись по коллекции
\pause
\usemintedstyle{lightbulb}
\begin{tcblisting}{listing engine=minted, minted language=java, minted options={fontsize=\scriptsize}, listing only, colback=codebg, colframe=codebg, rounded corners}
interface Iterator<T> {
  // Достигнут ли конец последовательности
  bool hasNext();
  // Возвращает следующий элемент последовательности
  T next();
}

while (it.hasNext()) {
  T value = it.next();
  ...
}
\end{tcblisting}
\end{itemize}
\end{frame}

\begin{frame}[fragile]
\frametitle{Итераторы: C++ style}
\begin{itemize}
\usemintedstyle{tango}
\item C++-style: итератор указывает на конкретный элемент, и позволяет итерироваться в разные стороны
\pause
\item Сам итератор не знает, где конец последовательности, нужны два итератора \mintinline{cpp}|begin/end| -- \textit{iterator range}
\pause
\item Имеют интерфейс, аналогичный указателям
\pause
\usemintedstyle{lightbulb}
\begin{tcblisting}{listing engine=minted, minted language=cpp, minted options={fontsize=\scriptsize}, listing only, colback=codebg, colframe=codebg, rounded corners}
class iterator {
public:
  // Возвращает текущий элемент последовательности
  T& operator*();

  // Переходит к следующему элементу последовательности
  iterator& operator++();
};

for (iterator it = begin; it != end; ++it) {
  T& value = *it;
  ...
}
\end{tcblisting}
\end{itemize}
\end{frame}

\begin{frame}[fragile]
\frametitle{Итераторы}
\begin{itemize}
\usemintedstyle{tango}
\item Все STL-контейнеры имеют свои итераторы (обычно это вложенные классы, напр. \mintinline{cpp}|std::vector<int>::iterator|)
\pause
\item Релизация итератора зависит от контейнера, но почти всегда это просто указатель на элемент / вершину списка / etc
\pause
\item \mintinline{cpp}|begin()/end()| -- получить итератор на начало/конец
\pause
\item \mintinline{cpp}|find(x)| -- возвращает итератор на найденный элемент (или на конец)
\pause
\item \mintinline{cpp}|erase(it)| -- удаляет элемент, на который указывает итератор
\pause
\item \mintinline{cpp}|insert(it, x)| -- вставляет элемент \textit{перед} тем, на который указывает итератор
\end{itemize}
\end{frame}

\begin{frame}[fragile]
\frametitle{End итератор}
\begin{itemize}
\usemintedstyle{tango}
\item `Итератор на конец' почти всегда указывает не на последний элемент, а \textit{за} последний элемент
\pause
\item Для \mintinline{cpp}|std::vector| -- на элемент \mintinline{cpp}|v.data() + v.size()|
\pause
\item Для \mintinline{cpp}|std::list| -- на фиктивную вершину
\pause
\item Это оказывается гораздо удобнее для многих алгоритмов
\pause
\item Например, если диапазон \mintinline{cpp}|[begin, end)| разбить на два некоторым элементом \mintinline{cpp}|middle|, то получатся два диапазона: \mintinline{cpp}|[begin, middle)| и \mintinline{cpp}|[middle, end)|
\end{itemize}
\end{frame}

\begin{frame}[fragile]
\frametitle{Инвалидация ссылок}
\begin{itemize}
\usemintedstyle{tango}
\item В чём проблема такого кода?
\usemintedstyle{lightbulb}
\begin{tcblisting}{listing engine=minted, minted language=cpp, minted options={fontsize=\scriptsize}, listing only, colback=codebg, colframe=codebg, rounded corners}
std::vector<int> v(10);
int& x = v[5];
v.push_back(1);
std::cout << x << std::endl;
\end{tcblisting}
\pause
\usemintedstyle{tango}
\item \mintinline{cpp}|push_back()| может перевыделить память \(\Longrightarrow\) ссылка \mintinline{cpp}|x| может указывать на уже освобождённую память
\pause
\item Говорят, что ссылка \textit{инвалидировалась}/\textit{инвалидирована}
\end{itemize}
\end{frame}

\begin{frame}[fragile]
\frametitle{Инвалидация итераторов}
\begin{itemize}
\usemintedstyle{tango}
\item В чём проблема такого кода?
\usemintedstyle{lightbulb}
\begin{tcblisting}{listing engine=minted, minted language=cpp, minted options={fontsize=\scriptsize}, listing only, colback=codebg, colframe=codebg, rounded corners}
std::vector<int> v(10);
// Итераторы вектора умеют прибавлять/вычитать
// числа, как указатели
auto it = v.begin() + 5;
v.push_back(1);
std::cout << *it << std::endl;
\end{tcblisting}
\pause
\usemintedstyle{tango}
\item \mintinline{cpp}|push_back()| может перевыделить память \(\Longrightarrow\) итератор \mintinline{cpp}|it| может указывать на уже освобождённую память
\pause
\item Говорят, что итератор \textit{инвалидировался}/\textit{инвалидирован}
\end{itemize}
\end{frame}

\begin{frame}[fragile]
\frametitle{Инвалидация итераторов}
\begin{itemize}
\usemintedstyle{tango}
\item Когда происходит инвалидация ссылок/итераторов?
\pause
\item Можно прочитать в документации\pause, но проще знать, как работают контейнеры :)
\pause
\item \mintinline{cpp}|std::vector| -- когда может перевыделиться память (при добавлении элемента), при удалении элемента (инвалидируются только итераторы \textit{после} удалённого элемента)
\pause
\item \mintinline{cpp}|std::deque| -- добавление/удаление инвалидирует итераторы (перевыделяется массив указателей на блоки), но не ссылки на элементы
\pause
\item \mintinline{cpp}|std::list| -- только если удаляется конкретный элемент (инвалидируется итератор на него, и только он)
\end{itemize}
\end{frame}

\begin{frame}[fragile]
\frametitle{Инвалидация итераторов}
\begin{itemize}
\usemintedstyle{tango}
\item В чём проблема такого кода?
\usemintedstyle{lightbulb}
\begin{tcblisting}{listing engine=minted, minted language=cpp, minted options={fontsize=\scriptsize}, listing only, colback=codebg, colframe=codebg, rounded corners}
// Удаляем все чётные элементы
for (auto it = v.begin(); it != v.end(); ++it)
  if (*it % 2 == 0)
    v.erase(it);
\end{tcblisting}
\pause
\usemintedstyle{tango}
\item \mintinline{cpp}|erase()| может инвалидировать итератор
\end{itemize}
\end{frame}

\begin{frame}[fragile]
\frametitle{Инвалидация итераторов}
\begin{itemize}
\usemintedstyle{tango}
\item \mintinline{cpp}|erase(it)| удаляет элемент, и возвращает итератор на следующий
\pause
\item В чём проблема такого кода?
\usemintedstyle{lightbulb}
\begin{tcblisting}{listing engine=minted, minted language=cpp, minted options={fontsize=\scriptsize}, listing only, colback=codebg, colframe=codebg, rounded corners}
// Удаляем все чётные элементы
for (auto it = v.begin(); it != v.end(); ++it) {
  if (*it % 2 == 0)
    it = v.erase(it);
}
\end{tcblisting}
\pause
\item Пропустим все элементы, следующие за чётными
\end{itemize}
\end{frame}

\begin{frame}[fragile]
\frametitle{Инвалидация итераторов}
\begin{itemize}
\usemintedstyle{tango}
\usemintedstyle{lightbulb}
\begin{tcblisting}{listing engine=minted, minted language=cpp, minted options={fontsize=\scriptsize}, listing only, colback=codebg, colframe=codebg, rounded corners}
// Удаляем все чётные элементы
for (auto it = v.begin(); it != v.end();) {
  if (*it % 2 == 0)
    it = v.erase(it);
  else
    ++it;
}
\end{tcblisting}
\end{itemize}
\end{frame}

\begin{frame}[fragile]
\frametitle{Инвалидация итераторов}
\begin{itemize}
\usemintedstyle{tango}
\item В чём проблема такого кода?
\usemintedstyle{lightbulb}
\begin{tcblisting}{listing engine=minted, minted language=cpp, minted options={fontsize=\scriptsize}, listing only, colback=codebg, colframe=codebg, rounded corners}
// Вставляем 0 перед каждым чётным элеменом
for (auto it = v.begin(); it != v.end(); ++it)
  if (*it % 2 == 0)
    v.insert(it, 0);
\end{tcblisting}
\pause
\usemintedstyle{tango}
\item \mintinline{cpp}|insert()| может инвалидировать итератор
\end{itemize}
\end{frame}

\begin{frame}[fragile]
\frametitle{Инвалидация итераторов}
\begin{itemize}
\usemintedstyle{tango}
\item \mintinline{cpp}|insert(it, value)| вставляет элемент, и возвращает новый итератор на него
\pause
\item В чём проблема такого кода?
\usemintedstyle{lightbulb}
\begin{tcblisting}{listing engine=minted, minted language=cpp, minted options={fontsize=\scriptsize}, listing only, colback=codebg, colframe=codebg, rounded corners}
// Вставляем 0 перед каждым чётным элеменом
for (auto it = v.begin(); it != v.end(); ++it)
  if (*it % 2 == 0)
    it = v.insert(it, 0);
\end{tcblisting}
\pause
\usemintedstyle{tango}
\item Бескончно будем вставлять 0 перед одним и тем же элементом
\end{itemize}
\end{frame}

\begin{frame}[fragile]
\frametitle{Инвалидация итераторов}
\begin{itemize}
\usemintedstyle{lightbulb}
\begin{tcblisting}{listing engine=minted, minted language=cpp, minted options={fontsize=\scriptsize}, listing only, colback=codebg, colframe=codebg, rounded corners}
// Вставляем 0 перед каждым чётным элеменом
for (auto it = v.begin(); it != v.end(); ++it) {
  if (*it % 2 == 0) {
    it = v.insert(it, 0);
    ++it;
  }
}
\end{tcblisting}
\end{itemize}
\end{frame}

\begin{frame}[fragile]
\frametitle{Ассоциативные контейнеры}
\begin{itemize}
\usemintedstyle{tango}
\item Хранят элементы по-особому для быстрого поиска
\pause
\item Упорядоченные:
\begin{itemize}
\item \mintinline{cpp}|set|, \mintinline{cpp}|multiset|, \mintinline{cpp}|map|, \mintinline{cpp}|multimap|
\pause
\item Реализуются в виде дерева (обычно -- красно-чёрное)
\pause
\item Требуют отношение порядка на элементах (\mintinline{cpp}|operator <|)
\pause
\item Операции выполняются за \(O(\log N)\)
\end{itemize}
\pause
\item Неупорядоченные:
\begin{itemize}
\item \mintinline{cpp}|unordered_set|, \mintinline{cpp}|unordered_multiset|, \mintinline{cpp}|unordered_map|, \mintinline{cpp}|unordered_multimap|
\pause
\item Реализуются в виде хеш-таблицы с закрытой адресацией
\pause
\item Требуют функцию, вычисляющую хеш
\pause
\item Операции выполняются за \(O(1)\) в среднем, за \(O(N)\) в худшем случае
\end{itemize}
\end{itemize}
\end{frame}

\begin{frame}[fragile]
\frametitle{Ассоциативные контейнеры}
\begin{itemize}
\usemintedstyle{tango}
\item \mintinline{cpp}|set|, \mintinline{cpp}|unordered_set| -- хранят уникальные элементы (без повторений)
\pause
\item \mintinline{cpp}|multiset|, \mintinline{cpp}|unordered_multiset| -- элементы могут повторяться
\pause
\item \mintinline{cpp}|map|, \mintinline{cpp}|unordered_map| -- хранят пары (ключ,значение) для уникальных ключей
\pause
\item \mintinline{cpp}|multimap|, \mintinline{cpp}|unordered_multimap| -- ключи могут повторяться
\end{itemize}
\end{frame}

\begin{frame}[fragile]
\frametitle{set}
\begin{itemize}
\usemintedstyle{lightbulb}
\begin{tcblisting}{listing engine=minted, minted language=cpp, minted options={fontsize=\scriptsize}, listing only, colback=codebg, colframe=codebg, rounded corners}
std::set<int> set;
set.insert(2);
set.insert(1);
set.insert(1); // вторая копия НЕ добавится
set.insert(9);
set.insert(5);

set.contains(1); // true
set.contains(3); // false
set.size(); // 4

// Выведется 1 2 5 9
for (auto it = set.begin(); it != set.end(); ++it) {
  std::cout << *it << " ";
}
\end{tcblisting}
\end{itemize}
\end{frame}

\begin{frame}[fragile]
\frametitle{multiset}
\begin{itemize}
\usemintedstyle{lightbulb}
\begin{tcblisting}{listing engine=minted, minted language=cpp, minted options={fontsize=\scriptsize}, listing only, colback=codebg, colframe=codebg, rounded corners}
std::multiset<int> set;
set.insert(2);
set.insert(1);
set.insert(1); // вторая копия добавится
set.insert(9);
set.insert(5);

set.contains(1); // true
set.contains(3); // false
set.size(); // 5

// Выведется 1 1 2 5 9
for (auto it = set.begin(); it != set.end(); ++it) {
  std::cout << *it << " ";
}
\end{tcblisting}
\end{itemize}
\end{frame}

\begin{frame}[fragile]
\frametitle{unordered\_set}
\begin{itemize}
\usemintedstyle{lightbulb}
\begin{tcblisting}{listing engine=minted, minted language=cpp, minted options={fontsize=\scriptsize}, listing only, colback=codebg, colframe=codebg, rounded corners}
std::unordered<int> set;
set.insert(2);
set.insert(1);
set.insert(1); // вторая копия НЕ добавится
set.insert(9);
set.insert(5);

set.contains(1); // true
set.contains(3); // false
set.size(); // 4

// Выведутся элементы в каком-то порядке,
// например 5 2 1 9
for (auto it = set.begin(); it != set.end(); ++it) {
  std::cout << *it << " ";
}
\end{tcblisting}
\end{itemize}
\end{frame}

\begin{frame}[fragile]
\frametitle{Константные итераторы}
\begin{itemize}
\usemintedstyle{tango}
\item У \mintinline{cpp}|*set| итераторы \textit{константные}, т.е. не позволяют модифицировать элемент, на который ссылаются
\pause
\item Иначе нарушились бы внутренние инварианты структуры данных (порядок в дереве, индекс в хеш-таблице)
\pause
\item У всех контейнеров есть отдельные типы \mintinline{cpp}|iterator| и \mintinline{cpp}|const_iterator| (у \mintinline{cpp}|*set| они совпадают)
\pause
\item \mintinline{cpp}|cbegin()/cend()| возвращают константные итераторы
\end{itemize}
\end{frame}

\begin{frame}[fragile]
\frametitle{map, std::pair}
\begin{itemize}
\usemintedstyle{tango}
\item \mintinline{cpp}|std::pair<X, Y>| -- вспомогательный класс, хранящий пару элементов
\pause
\usemintedstyle{lightbulb}
\begin{tcblisting}{listing engine=minted, minted language=cpp, minted options={fontsize=\scriptsize}, listing only, colback=codebg, colframe=codebg, rounded corners}
template <typename X, typename Y>
struct pair {
  X first;
  Y second;
};
\end{tcblisting}
\pause
\usemintedstyle{tango}
\item Различные варианты контейнера \mintinline{cpp}|map<Key, Value>| хранят в себе \mintinline{cpp}|std::pair<const Key, Value>|
\pause
\item В остальном \mintinline{cpp}|map| аналогичен \mintinline{cpp}|set|
\end{itemize}
\end{frame}

\begin{frame}[fragile]
\frametitle{map, std::pair}
\begin{itemize}
\usemintedstyle{lightbulb}
\begin{tcblisting}{listing engine=minted, minted language=cpp, minted options={fontsize=\scriptsize}, listing only, colback=codebg, colframe=codebg, rounded corners}
std::map<std::string, int> settings;

// Явное создание шаблона std::pair
settings.insert(std::pair<std::string, int>("volume", 100));

// Использование вспомогательной функции
settings.insert(std::make_pair("brightness", 80));

// Использование вывода параметров шаблонов (С++17)
settings.insert(std::pair("ui_scale", 3));

// Использование list initialization
settings.insert({"difficulty", 5});

// Итерация по элементам:
for (auto it = settings.begin(); it != settings.end(); ++it)
  std::cout << it->first << ": " << it->second << std::endl;
\end{tcblisting}
\end{itemize}
\end{frame}

\begin{frame}[fragile]
\frametitle{map: operator[]}
\begin{itemize}
\usemintedstyle{tango}
\item У \mintinline{cpp}|map| и \mintinline{cpp}|unordered_map| есть \mintinline{cpp}|operator[]|, упрощающий работу с ними:
\usemintedstyle{lightbulb}
\begin{tcblisting}{listing engine=minted, minted language=cpp, minted options={fontsize=\scriptsize}, listing only, colback=codebg, colframe=codebg, rounded corners}
std::map<std::string, int> settings;
settings["volume"] = 100;
settings["brightness"] = 80;

// Что произойдёт здесь?
std::cout << settings["ui_scale"];
\end{tcblisting}
\pause
\usemintedstyle{tango}
\item Обращение к несуществующему элементу через \mintinline{cpp}|operator[]| создаст его с помощью конструктора по умолчанию
\pause
\item \(\Longrightarrow\) Нужно использовать аккуратно, либо заменить на \mintinline{cpp}|map::find()|
\end{itemize}
\end{frame}

\begin{frame}[fragile]
\frametitle{Multiset}
\begin{itemize}
\usemintedstyle{tango}
\item У multi-версий ассоциативных контейнеров есть специальные методы для работы с последовательностью одинаковых элементов/ключей
\pause
\item \mintinline{cpp}|count(x)| -- сколько в контейнере содержится элементов, равных \mintinline{cpp}|x|
\pause
\item \mintinline{cpp}|equal_range(x)| -- вернуть пару \mintinline{cpp}|[begin, end)| итераторов на диапазон элементов, равных \mintinline{cpp}|x|
\pause
\item \textbf{NB:} у не-multi контейнеров эти методы тоже есть для единообразия интерфейса, но менее полезны (например, \mintinline{cpp}|count(x)| всегда возвращает 0 или 1)
\end{itemize}
\end{frame}

\end{document}